\begin{definition}[A locally oriented native checkerboard block]%
    \label{def:peps_kitaev_native_checkerboard_block}
    \lean{TNLean.PEPS.kitaevNativeBlockVertexEquiv,
        TNLean.PEPS.kitaevNativeCheckerboardSite,
        TNLean.PEPS.kitaevNativeBlockBoundary}
    \leanok
    \uses{def:peps_kitaev_checkerboard_tensor,def:peps_torus_plaquette_flux_geometry}
    Let the two torus periods be at least three. At an arbitrary site
    $v=(x,y)$, take the actual four-site plaquette with corners
    $(x,y)$, $(x+1,y)$, $(x+1,y+1)$ and $(x,y+1)$.
    Order the physical sites clockwise from the upper right.
    Place $T_1$ at the upper-right and lower-left corners, and $T_0$
    at the other two corners, with virtual legs in the order top, right,
    bottom, left. These are the local orientations in the blocking diagram
    of \cite[Section~7, ``Kitaev's code state'']{Schuch2010PEPS}.

    Read an actual crossing-bond assignment as four boundary pairs
    $N=(N_0,N_1)$, $E=(E_0,E_1)$, $S=(S_0,S_1)$ and $W=(W_0,W_1)$,
    ordered north, east, south, west, clockwise along each side.
    All coordinates are cyclic, so this definition also applies to a
    plaquette meeting either seam. The orientations are prescribed on this
    one block; no globally alternating tensor field on an odd torus is
    asserted.
\end{definition}

\begin{theorem}[Identification of the actual native checkerboard contraction]%
    \label{thm:peps_kitaev_native_checkerboard_identification}
    \lean{TNLean.PEPS.graphOpenRegionNetwork_kitaevNativeCheckerboardSite,
        TNLean.PEPS.graphOpenRegionNetwork_kitaevNativeCheckerboardSite_apply}
    \leanok
    \uses{def:peps_kitaev_native_checkerboard_block,
        def:peps_kitaev_checkerboard_block,
        thm:peps_kitaev_checkerboard_blocking,thm:peps_graph_open_region_contraction}
    For every actual crossing-bond assignment and every configuration of the
    four original physical spins, the open contraction on this native
    torus plaquette is exactly $B(N,E,S,W;\sigma)$ of
    Definition~\ref{def:peps_kitaev_checkerboard_block}. In particular,
    \begin{align}
        \mathcal A_v(N,E,S,W;\sigma)
        &=\delta_{N_1,N_0}\delta_{E_1,E_0}
          \delta_{S_1,S_0}\delta_{W_1,W_0}
          \delta_{\sigma,F(N_0,E_0,S_0,W_0)}.
    \end{align}
    This is an equality of actual open coefficients, with every original
    boundary label retained. It asserts neither a native boundary CNOT
    operator nor a globally tiled blocking identity.
\end{theorem}

\begin{proof}\leanok
    The four sites are distinct. Their internal edges are precisely the
    upper and lower horizontal bonds and the left and right vertical bonds.
    The period bounds exclude an additional internal bond across a seam.
    Split the incident labels into their internal and prescribed crossing
    labels, and enumerate the four internal bonds in clockwise order.
    The sixteen native leg reads give the four tensor factors in the
    definition of $B$. Reindexing the finite sum and product therefore gives
    the asserted equality. The coefficient formula follows from
    Theorem~\ref{thm:peps_kitaev_checkerboard_blocking}.
\end{proof}
